%% 第五章

\chapter{总结与展望}
\chaptermark{总结与展望}

\section{工作总结}

本文围绕低空无人机场景下多模态大模型的双认知时空推理能力评测与优化方法展开研究，整理并实现了 UAV-DualCog 双认知测试基准与对应能力优化方法。与传统通用视觉问答或视频理解任务不同，无人机在开放三维空间中的观测具有具身性、动态性和多视角性。模型不仅需要理解外部目标、地标和环境变化，还需要理解无人机自身相对位置、未来视角变化和飞行行为。因此，本文将无人机具身推理拆分为自我状态认知和环境情况认知两条主线，并在图像和视频两种媒介下开展统一评测。

在能力测试方面，本文构建了 UAV-DualCog 双认知测试基准，定义了四类图像任务和两类视频任务，形成“能力维度、观测媒介、证据形式”相结合的评测结构。数据构建上，本文实现了从 AerialVLN 仿真场景到测试样本的四阶段工具链，覆盖场景扫描与语义点云融合、地标挖掘与人工复核、行为驱动视频任务构建和证据感知图像任务构建，最终形成包含 12 个测试场景、512 个有效地标、4096 个图像任务样本和 2048 个视频任务样本的测试集。实验部分对闭源模型、开源模型和空间推理微调模型进行了统一评测，结果表明当前模型仍存在语义判断与显式证据脱节、自我状态认知弱于环境情况认知、图像与视频之间迁移不充分等问题。这些发现说明，面向无人机具身推理的评测不能只关注最终答案是否正确，还必须同时考察空间目标框、时间区间和双认知平衡。

在能力优化方面，本文根据评测发现构建了 UAV-DualCog-Train 图像任务训练集，并实现了以 Qwen 3.5 4B 为基座的两阶段优化路线。第一阶段采用真实图像输入的多模态 LoRA 监督微调，使模型学习稳定的结构化“答案+目标框”输出；第二阶段采用联合 GRPO，同时引入格式、语义和空间证据奖励，对语义-证据一致性进行进一步校正。任务级评测表明，优化模型在图像任务上取得显著直接收益，相比 Base 模型，图像选项准确率、BBox Acc@0.5 和 mIoU 分别提升 20.78、51.76 和 32.66 个百分点；在视频任务上，按复合行为识别、原子行为识别和地标可见性三类任务 Acc 算术平均计算的全局准确率提升 9.91 个百分点，但收益呈现迁移分化：自我行为识别与时间定位明显提升，环境可见性时序推理下降。该结果将第3章发现的能力短板转化为第4章中可执行、可验证的训练目标和实验证据，使论文形成从能力测试到能力优化的闭环。

\section{主要结论}

本文可以得到以下结论。

第一，双认知能力是无人机具身推理的必要组成部分。对于无人机智能体而言，仅理解外部环境并不足够，模型还必须理解自身状态以及自身运动如何改变观测内容。UAV-DualCog 通过自我状态认知和环境情况认知两条主线，将这一能力要求显式纳入评测。

第二，语义正确性不能替代证据可靠性。实验中大量模型能够给出正确选项，却不能提供准确目标框或时间区间。这说明多模态大模型在无人机场景中的失败方式并非简单“看不懂”，而是常表现为语义判断和空间、时间证据之间的不一致。

第三，当前模型在双认知能力上存在显著失衡。所有配对模型的环境情况认知总体准确率均高于自我状态认知，总体均值差达到 16.97 个百分点。这表明自我状态建模仍是当前无人机多模态推理中的核心短板。

第四，模型规模和已有空间推理微调不能自动解决无人机场景问题。实验结果显示，更大规模模型、MoE 模型或一般空间推理微调模型并不必然在 UAV-DualCog 中取得更好结果。面向无人机双认知和结构化证据输出的任务内优化仍具有必要性。

第五，能力优化应围绕双认知协同和证据一致性展开。评测结果已经说明，单纯提高答案准确率不足以支撑可靠无人机智能体。本文的优化实验进一步证明，证据感知图像训练能够显著提升图像侧答案-目标框一致性，并对视频自我行为时间定位产生正向迁移，尚不能直接改善环境可见性时序推理。这说明后续模型优化需要同时考虑图像空间证据、视频时间证据、自我状态认知和环境情况认知。

\section{研究局限}

本文仍存在需要进一步完善之处。第一，受训练性能、时间和算力条件限制，能力优化阶段优先覆盖图像任务，尚未直接引入视频任务样本进行联合训练。因此，优化模型虽然在图像任务和部分自我行为时间定位上取得收益，但在环境可见性时序推理中仍出现下降。第二，本文测试基准主要来源于 AerialVLN 仿真场景，能够保证数据构建可控和证据标注稳定，但与真实低空飞行中的光照、天气、传感器噪声和动态障碍物仍存在域差异。后续研究需要进一步引入真实飞行数据，检验双认知评测与能力优化方法在开放环境中的泛化能力。

\section{后续展望}

后续工作可从三个方向继续推进。

首先，可以进一步补充视频任务训练样本和跨媒介一致性约束。当前优化主要使用图像任务样本，已经能够提升图像空间证据质量并迁移到部分自我行为时间定位，但对环境可见性推理的帮助不足。后续可以引入视频片段、可见区间和跨帧目标一致性监督，使模型不仅能够在单帧中定位地标，也能够在连续视频中稳定跟踪地标并合并可见时间区间。

其次，可以围绕联合 GRPO 的奖励设计开展更细致的消融实验。本文当前采用格式、语义和目标框三项联合奖励，验证了整体路线的有效性，但不同奖励项对选项准确率、BBox 质量和视频迁移分化的具体贡献仍值得进一步拆解。后续可以分别比较仅格式奖励、语义奖励、空间证据奖励以及不同权重组合下的表现，从而更清楚地解释能力提升来源。

最后，可以扩大基座模型和场景来源，检验能力优化方法的泛化性。当前实验以 Qwen 3.5 4B 为主要优化基座，后续可在不同模型家族、Dense/MoE 架构和不同规模模型上复现实验，分析证据感知训练是否具有稳定迁移规律。同时，还可以逐步引入真实低空飞行数据和更复杂天气、光照、遮挡条件，使双认知评测与优化方法更接近真实无人机应用环境。
